\usepackage{wrapfig}
\usepackage{graphicx}
\usepackage{tcolorbox}
\usepackage{amsmath}
\usepackage{tabularx}

\newtcolorbox{insights}[1][]{%
    colback=blue!5!white, % Light blue background color of the box
    colframe=blue!75!black, % Darker blue border color of the box
    sharp corners, % Makes the corners of the box sharp
    toprule=0pt, % Remove top border
    bottomrule=0pt, % Remove bottom border
    leftrule=1.2pt, % Keep left border
    rightrule=1.2pt, % Keep right border
    #1 % Additional options
}

\setlength{\parskip}{2.5pt plus3pt minus2.5pt}
% 1.2pt：基础段落间距为 1.2 点（约 0.42 毫米）。
% plus3pt：允许 LaTeX 在排版需要时增加最多 3pt 的弹性间距（例如避免页面底部留白过多）。
% minus2.5pt：允许 LaTeX 在排版需要时减少最多 2.5pt 的弹性间距（例如避免内容被截断）。
% \setlength{\textfloatsep}{3pt plus2pt minus2pt} % 顶部/底部浮动体
% \setlength{\intextsep}{3pt plus2pt minus2pt}      % 中间浮动体
% 减少图表与上下文的间距（例如设为 10pt + 弹性值）

% 调整图片标题与内容的间距
% \setlength{\abovecaptionskip}{-10pt}  % 标题上方间距
% \setlength{\belowcaptionskip}{-10pt}   % 标题下方间距

% % 调整表格标题与内容的间距（表格标题通常在表格上方）
% \setlength{\abovecaptionskip}{-10pt}   % 表格标题上方间距
% \setlength{\belowcaptionskip}{-10pt}   % 表格标题下方间距